\documentclass[final,11pt,times]{elsarticle}

\usepackage{amssymb}
\usepackage{geometry}
\geometry{
 a4paper,
 left=23mm,
 right=23mm,
 top=10mm,
 bottom=15mm,
 headheight=60pt,
 includehead,includefoot,
 heightrounded,
}

\usepackage{fancyhdr}
\fancyhead{}
\pagestyle{fancy}
\fancyhead[c]{CIE53 Proceedings, October 20-23, 2026\\
Khalifa University, Abu Dhabi, UAE \\
}

\fancyhead[L]{\includegraphics[width=1.8cm]{logos/CIE_logo.png} \ \ \includegraphics[width=1.5cm]{logos/elsevier-official.pdf}}
\fancyhead[R]{\includegraphics[width=1.9cm]{logos/KULogo.png}}

\journal{CIE53 Proceedings}

\usepackage[colorlinks=true, pdfauthor={Djedjig, N.}]{hyperref}

\begin{document}

\begin{frontmatter}

\title{Reinforcement Learning for Adaptive Neighborhood Selection in Vehicle Routing with Synchronization Constraints}

\author[label1]{Nada Farah Djedjig}
\ead{nada.djedjig@example.com}

\author[label1]{Supervisor Name}
\ead{supervisor@example.com}

\author[label1]{Supervisor Name}
\ead{supervisor@example.com}

\address[label1]{Department of Computer Science, University Name, City, Country}

\begin{abstract}
The Vehicle Routing Problem with Time Windows and Synchronization constraints (VRPTW‑S) is central to home healthcare logistics. While metaheuristics like GVNS achieve strong performance, they rely on a fixed order of neighborhood operators, which cannot adapt to the changing search context. This paper presents the first systematic comparison of two reinforcement learning agents – tabular Q‑learning and Deep Q‑Network (DQN) – for dynamic operator selection in VRPTW‑S. We formulate the selection task as a Markov Decision Process, integrate both agents into a GVNS framework, and evaluate on Solomon‑based instances (50–400 customers). Results show that Q‑learning excels on small instances ($\leq 100$ customers) but fails to scale, while DQN remains stable and achieves consistent improvements up to 83\% on large instances (400 customers). Our findings provide practical guidance for selecting RL agents in adaptive routing systems for home healthcare.
\end{abstract}

\begin{keyword}
Vehicle Routing Problem \sep Time Windows \sep Synchronization \sep Reinforcement Learning \sep Q-learning \sep Deep Q-Network \sep Home Healthcare
\end{keyword}

\end{frontmatter}
\thispagestyle{fancy}

% ----------------------------------------------------------------------
% SECTION 1: INTRODUCTION
% ----------------------------------------------------------------------
\section{INTRODUCTION}
\label{sec:intro}

% Content to write:
% - Context: aging population, home healthcare expansion → VRPTW‑S as critical logistics problem.
% - Challenge: synchronization + time windows → NP‑hard; exact methods limited to ~20 customers.
% - Gap: state‑of‑the‑art metaheuristics (GVNS, ALNS) use static operator order; cannot adapt.
% - Contribution: first systematic comparison of Q‑learning vs DQN for operator selection in VRPTW‑S; MDP formulation; experiments on 50–400 customers.
% - Key insight: QL works for small problems, DQN scales.
% - Paper structure (brief).

% Source from report:
% - Chap. 1 §1.1 (HHC growth)
% - Chap. 1 §1.3 (complexity, exact methods limitations)
% - Chap. 4 §4.5 (static order limitation)
% - Chap. 5 intro (gap and contribution)

% Figures: none in Introduction.

% ----------------------------------------------------------------------
% SECTION 2: BACKGROUND
% ----------------------------------------------------------------------
\section{BACKGROUND}
\label{sec:background}

\subsection{VRPTW‑S Problem Definition}
% Write:
% - Graph model $G=(V,E)$; time windows $[e_i,l_i]$; service times $s_i$; synchronization constraints (strict, overlap, precedence).
% - MILP formulation with Big‑M (brief, not full derivation).
% - Complexity: NP‑hard; exact methods fail for realistic sizes.
% - Baseline metaheuristic: GVNS (shaking + VND + penalty self‑adjustment) – introduced in Section 4.

% Source:
% - Chap. 1 §1.2 and §1.3
% - Chap. 4 §4.2 (GVNS principle)

\subsection{Reinforcement Learning Fundamentals}
% Write:
% - Markov Decision Process (MDP) definition: $(S,A,P,R,\gamma)$.
% - Value‑based methods: Q‑learning (tabular), DQN (neural approximation).
% - Exploration vs exploitation: UCB, $\epsilon$-greedy.
% - Why RL is suitable for operator selection: sequential decision making under changing search states.

% Source:
% - Chap. 3 §3.2 (MDP)
% - Chap. 3 §3.4 (Q-learning, DQN)
% - Chap. 5 §5.3.2 (UCB, epsilon-greedy)

% Figures to include:
% - Figure 6 from report (agent-environment interaction loop) – label as Fig. 1 in this paper.

% ----------------------------------------------------------------------
% SECTION 3: MDP FORMULATION FOR OPERATOR SELECTION
% ----------------------------------------------------------------------
\section{MDP FORMULATION FOR OPERATOR SELECTION}
\label{sec:mdp}

% Write:
% - Define state space: 24‑dimensional feature vector with four categories (instance, solution quality, search progress, neighborhood success rates).
% - Define action space: 8 neighborhood operators (list them briefly).
% - Define reward function: normalized improvement $\Delta F_{\text{norm}}$, shaped with bonus/penalty.
% - Transition: deterministic.

% Source:
% - Chap. 5 §5.2.1 (state space, features)
% - Chap. 4 §4.3 (8 operators)
% - Chap. 5 §5.2.3 (reward)

% Figures: none (can add a small schematic if desired, but optional).

% ----------------------------------------------------------------------
% SECTION 4: METHODOLOGY – Q-LEARNING AND DQN AGENTS
% ----------------------------------------------------------------------
\section{METHODOLOGY: Q-LEARNING AND DQN AGENTS}
\label{sec:agents}

\subsection{Tabular Q-Learning}
% Write:
% - State discretization into 144 abstract states (feasibility flags, stagnation, effectiveness, instance size).
% - Q‑table structure, Bellman update rule.
% - Action selection: UCB (Upper Confidence Bound).
% - Hyperparameters: $\alpha=0.1$, $\gamma=0.95$, $c_{\text{UCB}}=1.0$.

% Source:
% - Chap. 5 §5.3.1 (discretization)
% - Chap. 5 §5.3.2 (UCB, hyperparameters)
% - Chap. 5 Table IX (hyperparameters)

% Figures to include:
% - Figure 20 from report (Q-learning internal architecture) – label as Fig. 2a.

\subsection{Deep Q-Network (DQN)}
% Write:
% - Neural network architecture: MLP with input (24 features) → two hidden layers (128 neurons, ReLU) → output (8 Q‑values).
% - Normalization (running mean/variance).
% - Experience replay buffer (size 50,000), target network (updated every 100 steps).
% - $\epsilon$-greedy exploration ($\epsilon$ decays from 1.0 to 0.05 over 500 steps).
% - Hyperparameters: learning rate $3\times10^{-4}$, batch size 32, Huber loss.

% Source:
% - Chap. 5 §5.4.1 (network architecture, normalization)
% - Chap. 5 §5.4.2 (replay, target network, epsilon)
% - Chap. 5 Table X (hyperparameters)

% Figures to include:
% - Figure from Chap. 5 showing MLP architecture (labeled "Figure : Architecture of the neural network...") – label as Fig. 2b.

\subsection{Integration into GVNS}
% Write:
% - At each VND step, the RL agent selects an operator.
% - The operator is applied, producing a new solution; reward and next state are computed.
% - The agent updates its Q‑table / network parameters.
% - The outer GVNS acceptance criterion decides whether to keep the new solution.

% Source:
% - Chap. 5 Fig. 19 (global architecture of RL‑driven GVNS)
% - Chap. 5 description between Fig. 18 and Fig. 19.

% Figures to include:
% - Figure 19 from report (RL‑driven GVNS architecture) – label as Fig. 3.

% ----------------------------------------------------------------------
% SECTION 5: EXPERIMENTAL RESULTS
% ----------------------------------------------------------------------
\section{EXPERIMENTAL RESULTS}
\label{sec:results}

\subsection{Experimental Setup}
% Write:
% - Instances: Solomon (C, R, RC) with 50, 100, 200, 400 customers; synchronization added (5/10/20/40 customer pairs).
% - Baseline: classical GVNS with fixed operator order.
% - Each instance run 5 times; metrics: % gap relative to baseline, success rate (proportion of instances where RL improves).

% Source:
% - Chap. 4 §4.6 (instance generation)
% - Chap. 5 §5.3.2 and §5.4.2 (experimental protocol)

% Figures/Tables: none in this subsection.

\subsection{Q-Learning Results}
% Write summary (textual, with reference to table):
% - 50 customers: outperforms GVNS in ~72% of instances; gaps -0.1% to -2%.
% - 100 customers: ~55% improvement; more variable, worst gap +5.5%.
% - 200–400 customers: fails to scale; worst gap +24.8% on 400.

% Source:
% - Chap. 5 §5.3.2 (tables mentioned for 50, 100, 200, 400)

% Table to include:
% - Create a new table (Table 1) summarizing QL results: columns [Instance size, % better than GVNS, Best gap (%), Worst gap (%)].
%   Data from Chap. 5 §5.3.2.

\subsection{DQN Results}
% Write summary:
% - 50 customers: ~45% improvement (less consistent than QL).
% - 100 customers: ~67% improvement; gaps up to -2.5%, stable.
% - 200–400 customers: ~83% improvement; gaps -0.5% to -1.8%, reliable.

% Source:
% - Chap. 5 §5.4.2 (tables for 50,100,200,400)

% Table to include:
% - Table 2 (same structure as Table 1, for DQN).

\subsection{Comparison and Discussion}
% Write:
% - Compare QL vs DQN: QL is simpler and better for small problems but collapses due to state explosion; DQN generalizes via neural approximation and scales.
% - Runtime observations: DQN often finds good solutions faster than GVNS.
% - Practical implication: For real‑world large‑scale HHC, DQN‑based adaptive selection is preferred.

% Source:
% - Chap. 5 conclusion synthesis.

% Figures to include:
% - Bar chart or boxplot summarizing % improvement vs instance size (you can create from data in Chap. 5). 
% - If not yet created, mention "Figure 4 will show performance comparison."

% ----------------------------------------------------------------------
% SECTION 6: CONCLUSION
% ----------------------------------------------------------------------
\section{CONCLUSION}
\label{sec:conclusion}

% Write:
% - Recap: first systematic comparison of QL and DQN for operator selection in VRPTW‑S.
% - QL works well for small instances (≤100), DQN scales to 400+ customers with consistent gains.
% - Future work: richer state representations (graph neural networks), multi‑objective optimisation, integration with other learning components (hyperparameter control).
% - Link to CIE53 theme: adaptive routing reduces travel and fuel consumption → sustainable healthcare logistics.

% Source:
% - Chap. 5 conclusion; Chap. 8 (future perspectives).

% ----------------------------------------------------------------------
% ACKNOWLEDGMENTS
% ----------------------------------------------------------------------
\section*{Acknowledgments}
% Optional – thank supervisor, funding, etc.

% ----------------------------------------------------------------------
% REFERENCES
% ----------------------------------------------------------------------
% Use the bibliography style from the template (elsarticle-num-names).
% Provide a .bib file or write \bibitem entries.
% At minimum, cite:
% - Bredström & Rönnqvist (2008) – foundational VRPTW‑S
% - Masmoudi et al. (2025) – GVNS baseline
% - Liu et al. (2019) – ALNS for VRPTW‑S
% - Sutton & Barto (2018) – RL textbook
% - Mnih et al. (2015) – DQN Nature paper
% - Bello et al. (2017) – Pointer Networks
% - Kool et al. (2019) – Attention Model
% - Solomon (1987) – instances
% - Afifi et al. (2016) – SA for VRPTW‑S

% The following is a placeholder – replace with actual \bibitem or .bib file.

\section*{References}

\begin{thebibliography}{99}

\bibitem{bre08} D. Bredstr\"om and M. R\"onnqvist, ``Combined vehicle routing and scheduling with temporal precedence and synchronization constraints,'' \textit{European Journal of Operational Research}, vol. 191, no. 1, pp. 19–31, 2008.

\bibitem{mas25} M. Masmoudi, R. Borchani, and B. Jarboui, ``Generalized variable neighborhood search algorithm for vehicle routing problem with time windows and synchronization,'' \textit{Computers \& Operations Research}, vol. 183, 107193, 2025.

\bibitem{liu19} R. Liu, Y. Tao, and X. Xie, ``An adaptive large neighborhood search heuristic for the vehicle routing problem with time windows and synchronized visits,'' \textit{Computers \& Operations Research}, vol. 101, pp. 250–262, 2019.

\bibitem{sut18} R. S. Sutton and A. G. Barto, \textit{Reinforcement Learning: An Introduction}, 2nd ed. MIT Press, 2018.

\bibitem{mnih15} V. Mnih et al., ``Human-level control through deep reinforcement learning,'' \textit{Nature}, vol. 518, no. 7540, pp. 529–533, 2015.

\bibitem{bel17} I. Bello, H. Pham, Q. V. Le, M. Norouzi, and S. Bengio, ``Neural combinatorial optimization with reinforcement learning,'' in \textit{ICLR}, 2017.

\bibitem{kool19} W. Kool, H. van Hoof, and M. Welling, ``Attention, learn to solve routing problems!'' in \textit{ICLR}, 2019.

\bibitem{solo87} M. M. Solomon, ``Algorithms for the vehicle routing and scheduling problems with time window constraints,'' \textit{Operations Research}, vol. 35, no. 2, pp. 254–265, 1987.

\bibitem{afifi16} S. Afifi, D.-C. Dang, and A. Moukrim, ``Heuristic solutions for the vehicle routing problem with time windows and synchronized visits,'' \textit{Optimization Letters}, vol. 10, no. 3, pp. 511–525, 2016.

\end{thebibliography}

\end{document}